\RequirePackage{pdfmanagement}
\SetKeys[document/metadata]{pdfversion=1.7, pdfstandard=A-2u, lang=en-GB}
\input{preamble/pdfA-compliance/glyphtounicode}
\documentclass[11pt,a4paper]{report}

\newcommand*{\PaperTitle}{\(n\) cooled CSTRs in series: \(3^n\) steady states, \(2^n\) of them stable}
\newcommand*{\PaperTitlePlain}{n cooled CSTRs in series: 3^n steady states, 2^n of them stable}
\newcommand*{\PaperShortTitle}{\(3^n\) steady states in \(n\) CSTRs in series}
\newcommand*{\PaperSubtitle}{chemistry-open-problems \textperiodcentered{} 5 November 2026 \textperiodcentered{} solved (preliminary, verification pending)}
\newcommand*{\PaperCredits}{Solver: Claude Opus 5.5 (effort xhigh) \textperiodcentered{} Verification: pending}
\newcommand*{\PaperKeywords}{continuous stirred tank reactor, CSTR cascade, steady-state multiplicity, stability, Arrhenius kinetics, interval arithmetic}
\input{chemop-preamble}
\addbibresource{references.bib}

\newcommand*{\Da}{\mathrm{Da}_0}
\newcommand*{\eps}{\varepsilon}

\begin{document}
\makepapertitle

\begin{paperabstract}
In his 1987 thesis on two continuous stirred tank reactors (CSTRs) in series with a first-order exothermic
reaction, each with its own cooling jacket, B.~Y. Bisbis showed that two reactors can have nine steady states and
conjectured that \(n\) reactors have at most \(3^n\) steady states, at most \(2^n\) of them stable. We prove the
conjecture for the model of the thesis, with Arrhenius kinetics and \(\eps=RT_0/E>0\). The upper bounds hold for all
parameter values, and also in the exponential limit \(\eps=0\): one reactor has at most three steady states for any
feed, the Jacobian of the cascade is block lower triangular, and the middle state of a reactor with three states is a
saddle. The bounds are attained for every \(n\) and every \(\eps>0\). For \(n\) identical reactors (equal volumes, equal
jackets, equal coolant temperature) we give explicit parameter values with exactly \(3^n\) steady states, all
hyperbolic, of which exactly \(2^n\) are asymptotically stable. The construction uses the bounded Arrhenius factor.
The Damköhler number at infinite temperature is \(k_0V/F=1/(4n)\), so a reactor converts less than \(1/(4n+1)\) of
the reactant it receives, even when ignited, and at least \(e^{-1/4}\) of the feed reaches the last reactor. Strong
cooling makes each reactor almost insensitive to the temperature of its feed. The proof reduces to five
elementary inequalities, checked in ball arithmetic. Interval certificates (Krawczyk test on the original equations)
confirm the \(3^n\) states and the \(2^n\) stable ones for \(n\le 6\). Two certified one-reactor examples show that the
outer states need not be stable: three states with an unstable upper focus, and three states that are all unstable.
The parameters that attain \(3^n\) are far from practical operation: ignited reactors run above \(3E/R\). The
exponential limit \(\eps=0\), which the thesis uses in its numerical example, is not covered by the construction.
\end{paperabstract}

\section*{The question}

B.~Y. Bisbis, \emph{On the steady states and dynamic behavior of two CSTR's in series}, PhD thesis, University of
Missouri--Columbia, December 1987, supervisor P.~C.-H. Chan~\autocite{Bisbis1987}. We read the thesis in full (open
access). Chapter~II proves, with the sequential singularity theory of Dangelmayr and
Stewart~\autocite{DangelmayrStewart1984,DangelmayrStewart1985}, that two reactors can have nine steady states, and gives
a numerical example with \(\eps=0\). The concluding remarks of Chapter~II (p.~27) state:
\begin{openquote}
``For the case of n - CSTR's in series, one may conjecture that reiteration of the same methods used here, along with
an induction reasoning, would lead to the following : The maximum possible number of steady states is \(3^n\), of which
at most only \(2^n\) are stable, the rest unstable.''
\end{openquote}
The abstract repeats that ``some global statements about maximum multiplicity in the case of n - CSTR's in series are
made''. On the parameters (p.~29): ``we have the same reaction taking place in both reactors and therefore kinetic
parameters such as B, \(\Da\) and \(\eps\) are chosen only once. [\ldots] The rest of the parameters \(\alpha\),
\(\eta_1\), \(\eta_2\), \(\beta_1\) and \(\beta_2\) may be chosen independently''. The residence time \(\tau\) is also
common, since one flow rate passes through all reactors.

\paragraph{The model.} The thesis writes the balances of reactor \(i\) with volume \(V_i\), jacket coefficient
\(u_ia_i\) and coolant temperature \(T_{ci}\) (eqs.~(1)--(4)) and the dimensionless form (eqs.~(5)--(8), p.~4; the scaling of Uppal, Ray and Poore~\autocite{UppalRayPoore1976}) with
\(\theta_i=\frac{E}{RT_0}\frac{T_i-T_0}{T_0}\), \(y_i=(c_{A0}-c_{Ai})/c_{A0}\), \(\eps=RT_0/E\),
\(B=(-\Delta H)c_{A0}/(\eps\rho C_pT_0)\), \(\eta_i=\frac{E}{RT_0}\frac{T_{ci}-T_0}{T_0}\), \(\Da=k_0e^{-1/\eps}t_r\),
\(\tau=V_1/(Ft_r)\), \(\alpha=V_2/V_1\). The \(n\)-reactor version is, for \(k=1,\dots,n\),
\begin{equation}\label{eq:model}
\begin{aligned}
\frac{dy_k}{dt}&=\frac{y_{k-1}-y_k}{\alpha_k}+\Da\tau(1-y_k)E(\theta_k),\\
\frac{d\theta_k}{dt}&=\frac{\theta_{k-1}-\theta_k}{\alpha_k}+\Da\tau B(1-y_k)E(\theta_k)
-\frac{\beta_k\tau}{\alpha_k}(\theta_k-\eta_k),
\end{aligned}
\qquad E(\theta)=\exp\Bigl(\frac{\theta}{1+\eps\theta}\Bigr),
\end{equation}
with \(y_0=\theta_0=0\) (fresh feed at \(T_0\)), \(\alpha_1=1\) and \(\alpha_k=V_k/V_1\). For \(n=2\) this is eqs.~(5)--(8)
of the thesis verbatim. The parameters are \(\eps>0\), \(B>0\), \(\Da>0\), \(\tau>0\) (common), and \(\alpha_k>0\),
\(\beta_k\ge0\), \(\eta_k>-1/\eps\) (per reactor; the last condition says \(T_{ck}>0\)). The phase space is
\(\{(y_k,\theta_k)_{k\le n}:\ 1+\eps\theta_k>0\}\), that is, positive absolute temperatures. A steady state is an
equilibrium of~\eqref{eq:model}. \emph{Stable} means stable in the sense of Lyapunov.
For the upper bounds we also allow \(\eps=0\), \(E(\theta)=e^\theta\), \(\theta_k,\eta_k\in\mathbb R\) (the
Frank-Kamenetskii approximation, used in the numerical example of the thesis).

\begin{remark}[note={on \(\beta_k\)}]
The thesis lists \(\beta_i=u_ia_it_r/(V_i\rho C_p)\), whereas eq.~(8) as printed corresponds to normalising by
\(V_1\). With the other convention, \(\beta_k\tau/\alpha_k\) becomes \(\beta_k\tau\) in~\eqref{eq:model}. Nothing below
depends on this: in \zcref{thm:upper} only the number \(b_k\ge0\) defined in~\eqref{eq:reduced} matters (it is
\(\alpha_k\beta_k\tau\) with the other convention), and the construction of \zcref{thm:family} has \(\alpha_k=1\), where
both conventions agree.
\end{remark}

\paragraph{Prior work.} Varma~\autocite{Varma1980} gave the bound \(2^{n+1}-1\) steady states, at most \(n+1\) stable,
for \(n\) reactors in series, ``for the adiabatic and the isothermal occurrence of the reaction'' according
to~\autocite{Bisbis1987} (see also~\autocite{GalanicsSzalai2024}). In those cases each reactor has a single state
variable. We have not read~\autocite{Varma1980}, which is paywalled. Svoronos, Aris and
Stephanopoulos~\autocite{Svoronos1982} studied two identical tanks in the infinite-activation-energy approximation
(\(\eps=0\), according to~\autocite{Bisbis1987}); their abstract reports up to seven steady states for a given
Damköhler number, at most three of them stable. Kubíček et al.~\autocite{Kubicek1980} studied two nonadiabatic
reactors numerically. Bisbis~\autocite{Bisbis1987} and Retzloff, Chan, Chicone and Papick~\autocite{Retzloff1992} proved that
nine states occur for two jacketed reactors. Li~\autocite{Li1994} treats isothermal reactions in two CSTRs; his
abstract (paraphrased in secondary sources) extends the count to \(3^N\) states, \(2^N\) stable, for \(N\) isothermal
reactors. That is a different model, without energy balance. We found no proof, for any \(n\ge3\), that \(3^n\)
states are attained in the jacketed nonisothermal model.

\begin{theorem}[note={main result}]\label{thm:main}
\begin{enumerate}
\item For every \(n\ge1\) and all parameter values (\(\eps\ge0\)), system~\eqref{eq:model} has at most \(3^n\) steady
states, and at most \(2^n\) of them are stable (\zcref{thm:upper}).
\item For every \(n\ge1\) and every \(\eps>0\), put \(m=\max\{1,\eps\}\) and
\begin{gather*}
B=\frac{6400\,n\,m}{\eps^2},\qquad \Da=\frac{e^{-1/\eps}}{4n},\qquad \tau=1,\\
\alpha_k=1,\qquad \beta_k=200m-1,\qquad \eta_k=\frac{1-20\eps}{20\eps^2}\qquad(k=1,\dots,n).
\end{gather*}
Then~\eqref{eq:model} has exactly \(3^n\) steady states. All are hyperbolic; exactly \(2^n\) are asymptotically
stable, and the other \(3^n-2^n\) are unstable, each with a positive real eigenvalue (\zcref{thm:family}).
\item Outer steady states need not be stable: there are parameters with \(n=1\) and three steady states of which only
one is stable, and parameters with three steady states none of which is stable (\zcref{prop:unstable}).
\end{enumerate}
\end{theorem}
In physical units, the family of part~2 has identical reactors with \(k_0V/F=1/(4n)\), \(UA/(F\rho C_p)=200m-1\),
coolant temperature \(T_c=E/(20R)\) and adiabatic temperature rise \((-\Delta H)c_{A0}/(\rho C_p)=6400\,n\,m\,E/R\). The
feed temperature \(T_0=\eps E/R\) is arbitrary. For \(\eps=1/20\) the coolant is at the feed temperature
(\(\eta_k=0\)).

\section{One reactor}\label{sec:one}

Fix a reactor \(k\) and its inlet \((y_{k-1},\theta_{k-1})\), and write \(c_{k-1}=1-y_{k-1}\). Set
\begin{equation}\label{eq:reduced}
D_k=\alpha_k\Da\tau,\quad b_k=\beta_k\tau,\quad \lambda_k=\frac1{1+b_k}\in(0,1],\quad
g_k(\theta)=\frac{D_kE(\theta)}{1+D_kE(\theta)},
\end{equation}
\begin{equation}\label{eq:phi}
\varphi_k(\theta)=\lambda_k\theta_{k-1}+(1-\lambda_k)\eta_k+\lambda_kBc_{k-1}\,g_k(\theta)-\theta .
\end{equation}

\begin{lemma}[note={reduction}]\label{lem:reduction}
Let \(c_{k-1}>0\). The pair \((y_k,\theta_k)\) makes the right-hand side of reactor \(k\) in~\eqref{eq:model} vanish
if and only if \(\varphi_k(\theta_k)=0\) and \(1-y_k=c_{k-1}/(1+D_kE(\theta_k))\). In particular \(0<1-y_k\le c_{k-1}\)
at a steady state, and the steady states of~\eqref{eq:model} are the sequences obtained by choosing, for
\(k=1,\dots,n\) in turn, a zero \(\theta_k\) of \(\varphi_k\) (which depends on the choices for reactors \(1,\dots,k-1\)).
\end{lemma}
\begin{proof}
Multiply the first equation by \(\alpha_k\): it vanishes iff \(c_{k-1}-c_k=D_kc_kE(\theta_k)\) with \(c_k=1-y_k\), that
is \(c_k=c_{k-1}/(1+D_kE)\). On this curve, \(\alpha_k\) times the second equation equals \((1+b_k)\varphi_k(\theta_k)\)
(check A2 of \texttt{check\_model\_algebra.py}). Since \(c_0=1\), induction gives \(c_k>0\) for all \(k\).
\end{proof}

\begin{lemma}[note={at most three states}]\label{lem:three}
Let \(B>0\), \(c_{k-1}>0\), \(\eps\ge0\). Then \(\varphi_k\) has at most three zeros on \(I=\{\theta:1+\eps\theta>0\}\).
If it has three zeros \(\theta^1<\theta^2<\theta^3\), then \(\varphi_k'(\theta^1)<0\), \(\varphi_k'(\theta^2)>0\) and
\(\varphi_k'(\theta^3)<0\).
\end{lemma}
\begin{proof}
Write \(g=g_k\). A direct computation (checks A4) gives
\[
g'=\frac{g(1-g)}{(1+\eps\theta)^2}>0,\qquad
g''=\frac{g(1-g)}{(1+\eps\theta)^4}\,s(\theta),\qquad s(\theta)=1-2g(\theta)-2\eps(1+\eps\theta).
\]
The function \(s\) is strictly decreasing on \(I\), because \(g\) is strictly increasing and \(\eps(1+\eps\theta)\) is
nondecreasing. So \(g''>0\) on an interval \(I_-\) (possibly empty) and \(g''<0\) on the complementary interval \(I_+\)
to its right, up to at most one point. Since \(\varphi_k''=\lambda_kBc_{k-1}g''\), the derivative \(\varphi_k'\) is
strictly increasing on \(I_-\) and strictly decreasing on \(I_+\). It has at most two zeros, and by Rolle's theorem
\(\varphi_k\) has at most three.

Suppose there are three zeros. By Rolle, \(\varphi_k'\) vanishes at some \(\xi_1\in(\theta^1,\theta^2)\) and some
\(\xi_2\in(\theta^2,\theta^3)\). These are its only zeros, so \(\varphi_k'(\theta^j)\ne0\) and \(\varphi_k\) changes
sign at each \(\theta^j\). Since \(g<1\), \(\varphi_k(\theta)\to-\infty\) as \(\theta\to\infty\). Hence
\(\varphi_k<0\) on \((\theta^3,\infty)\), \(\varphi_k>0\) on \((\theta^2,\theta^3)\) and \(\varphi_k<0\) on
\((\theta^1,\theta^2)\), and the signs of \(\varphi_k'(\theta^j)\) follow.
\end{proof}

\begin{lemma}[note={linearisation}]\label{lem:jac}
Order the variables as \((y_1,\theta_1,\dots,y_n,\theta_n)\). The Jacobian of~\eqref{eq:model} is block lower
triangular with \(2\times2\) diagonal blocks \(J_k\); the block of reactor \(k\) depends only on
\((y_k,\theta_k)\) and the parameters. At a steady state, with \(\mathcal E_k=D_kE(\theta_k)\),
\begin{align*}
\alpha_k^2\det J_k&=-(1+\mathcal E_k)(1+b_k)\,\varphi_k'(\theta_k),\\
\alpha_k\operatorname{tr}J_k&=(1+\mathcal E_k)(1+b_k)\Bigl(\varphi_k'(\theta_k)-\lambda_k+\frac{\mathcal E_k}{1+\mathcal E_k}\Bigr).
\end{align*}
Hence the eigenvalues of the Jacobian are those of the blocks, and
\begin{enumerate}
\item if \(\varphi_k'(\theta_k)>0\), then \(J_k\) has a positive and a negative real eigenvalue;
\item \(J_k\) is Hurwitz (both eigenvalues with negative real part) if and only if
\(\varphi_k'(\theta_k)<\min\{0,\ \lambda_k-\mathcal E_k/(1+\mathcal E_k)\}\).
\end{enumerate}
\end{lemma}
\begin{proof}
Reactor \(k\) involves only \((y_{k-1},\theta_{k-1},y_k,\theta_k)\), which gives the block structure (check A7 for
\(n=3\)). The two identities are checks A3. A real \(2\times2\) matrix is Hurwitz iff its determinant is positive and
its trace negative, and has real eigenvalues of opposite signs iff its determinant is negative.
\end{proof}

\section{Upper bounds}\label{sec:upper}

\begin{theorem}[note={upper bounds}]\label{thm:upper}
Let \(n\ge1\), \(\eps\ge0\), \(B,\Da,\tau,\alpha_k>0\), \(\beta_k\ge0\), \(\eta_k\in\mathbb R\).
\begin{enumerate}
\item System~\eqref{eq:model} has at most \(3^n\) steady states.
\item A steady state at which \(\varphi_k'(\theta_k)>0\) for some \(k\) has a positive real eigenvalue, so it is
unstable. At most \(2^n\) steady states satisfy \(\varphi_k'(\theta_k)\le0\) for all \(k\). Hence at most \(2^n\) steady
states are stable.
\item If there are exactly \(3^n\) steady states, then exactly \(3^n-2^n\) of them have a reactor with
\(\varphi_k'(\theta_k)>0\), and these are unstable.
\end{enumerate}
\end{theorem}
\begin{proof}
By \zcref{lem:reduction}, the steady states form a rooted tree of depth \(n\). A node at depth \(k-1\) is a steady
state of reactors \(1,\dots,k-1\); its children are the zeros of the corresponding \(\varphi_k\). By
\zcref{lem:three} every node has at most three children, so there are at most \(3^n\) leaves. This proves part~1.

If \(\varphi_k'(\theta_k)>0\), then \(J_k\) has a positive eigenvalue (\zcref{lem:jac}), hence so has the full
Jacobian, and a steady state at which the Jacobian has an eigenvalue with positive real part is unstable (Lyapunov;
see~\autocite[§2.9]{Perko2001}). Every node has at most two children with \(\varphi_k'\le0\): if it has three
children, the middle one has \(\varphi_k'>0\) by \zcref{lem:three}. So at most \(2^n\) leaves have
\(\varphi_k'\le0\) along the whole path, which proves part~2. If there are \(3^n\) leaves, every node has exactly three
children, exactly two of them with \(\varphi_k'<0\), and exactly \(2^n\) paths avoid the middle child; this proves
part~3.
\end{proof}

\begin{remark}
The upper bound uses only that each reactor has at most three steady states for every feed, and the triangular
structure. The same proof applies if the reactors differ in any per-reactor parameter, and for \(\eps=0\).
\end{remark}

\begin{proposition}[note={outer states can be unstable}]\label{prop:unstable}
Take \(n=1\), \(\eps=1/20\), \(\tau=1\).
\begin{enumerate}
\item For \(B=20\), \(\Da=1/10\), \(\beta_1=2\), \(\eta_1=-1\) there are exactly three steady states,
\(\theta\approx-0.1279\), \(2.3968\), \(4.9257\): a hyperbolic sink, a saddle, and an unstable focus (trace
\(\approx1.596\), determinant \(\approx7.815\)).
\item For \(B=30\), \(\Da=69/500\), \(\beta_1=5\), \(\eta_1=-7/20\) there are exactly three steady states,
\(\theta\approx1.3456\), \(1.7744\), \(2.6969\): an unstable node, a saddle, and an unstable node. No steady state is
stable.
\end{enumerate}
\end{proposition}
\begin{proof}
Computer-assisted (\texttt{certify\_unstable\_outer\_state.py}). For each example, ball arithmetic shows three sign
changes of \(\varphi_1\) on a grid of step \(1/32\), so by \zcref{lem:three} there are exactly three zeros. Each steady
state is then verified by a Krawczyk test in \((y,\theta)\), and the trace, determinant and discriminant of \(J_1\) are
enclosed on the verified box; their signs give the stated types. The enclosures are in
\texttt{certificates/unstable\_outer\_state.json}.
\end{proof}

So the two outer states of a reactor can both be unstable. The bound \(2^n\) is therefore not reached merely because
there are \(3^n\) steady states. \zcref{thm:family} shows that it is reached.

\section{Attainment for every \texorpdfstring{\(n\)}{n}}\label{sec:attain}

From now on \(\eps>0\). We use the reduced temperature \(z=RT/E=\eps(1+\eps\theta)>0\), so that
\(E(\theta)=e^{1/\eps}e^{-1/z}\), and the reduced concentration \(c=1-y\). For reactor \(k\) set
\[
K_k=D_ke^{1/\eps}=\frac{k_0V_k}F,\quad \Delta=\eps^2B=\frac{R\,(-\Delta H)c_{A0}}{E\rho C_p},\quad
z_{c,k}=\eps(1+\eps\eta_k)=\frac{RT_{ck}}E>0 .
\]
\(K_k\) is the Damköhler number at infinite temperature and \(\Delta\) the adiabatic temperature rise in units of
\(E/R\). The inlet of reactor \(k\) is \((c_{k-1},z_{k-1})\), with \((c_0,z_0)=(1,\eps)\).

\begin{lemma}[note={reduced variables}]\label{lem:zform}
With \(k(z)=K_ke^{-1/z}\) and
\[
F_k(z)=\lambda_kz_{k-1}+(1-\lambda_k)z_{c,k}+\lambda_k\Delta c_{k-1}\,\frac{k(z)}{1+k(z)}-z ,
\]
one has \(F_k(z)=\eps^2\varphi_k(\theta)\) and \(F_k'(z)=\varphi_k'(\theta)\) for \(z=\eps(1+\eps\theta)\). The outlet
of a steady state of reactor \(k\) is \(c_k=c_{k-1}/(1+k(z_k))\). Moreover, with \(X_k=\lambda_k\Delta K_k\) and
\(\psi(z)=e^{-1/z}/z^2\),
\[
F_k'(z)=X_kc_{k-1}\frac{\psi(z)}{(1+k(z))^2}-1 .
\]
\end{lemma}
\begin{proof}
Substitution (checks A1 and A5).
\end{proof}

Note that \(\psi'(z)=e^{-1/z}z^{-4}(1-2z)\) (check A6). So \(\psi\) increases on \((0,1/2]\) and decreases on
\([1/2,\infty)\). Also \(k/(1+k)<K_k/(1+K_k)\) and \(e^{-1/z}/(1+K_ke^{-1/z})=1/(e^{1/z}+K_k)<1/(1+K_k)\) for all
\(z>0\).

\begin{lemma}[note={window lemma}]\label{lem:window}
Fix a reactor with parameters \(\lambda\in(0,1)\), \(K>0\), \(z_c>0\), \(\Delta>0\), and put \(X=\lambda\Delta K\).
Let \(0<a<1/2<b\), \(0<c_-\le1\) and \(\zeta_+>0\) satisfy
\begin{align}
&\lambda\zeta_++z_c+Xe^{-1/a}<a, \tag{W1}\label{W1}\\
&Xc_-e^{-1/b}>(1+K)\,b, \tag{W2}\label{W2}\\
&(1+K)X\psi(a)<1\quad\text{and}\quad(1+K)X\psi(b)<1. \tag{W3}\label{W3}
\end{align}
Then for every inlet \((c,\zeta)\) with \(c_-\le c\le1\) and \(0<\zeta\le\zeta_+\):
\begin{enumerate}
\item the reactor has exactly three steady states, with temperatures \(z_L\in(0,a)\), \(z_M\in(a,b)\) and
\(z_H\in(b,\infty)\);
\item at \(z_L\) and \(z_H\) the block \(J\) is Hurwitz; at \(z_M\) it has real eigenvalues of opposite signs;
\item every outlet satisfies \(c_-/(1+K)\le c/(1+K)<c_{\mathrm{out}}<c\le1\) and
\(0<z_{\mathrm{out}}<\lambda\zeta_++(1-\lambda)z_c+X/(1+K)\).
\end{enumerate}
\end{lemma}
\begin{proof}
Write \(F(z)=\lambda\zeta+(1-\lambda)z_c+Xc\,e^{-1/z}/(1+Ke^{-1/z})-z\), which is \(F_k\) of \zcref{lem:zform}
(since \(\lambda\Delta c\,k/(1+k)=Xc\,e^{-1/z}/(1+Ke^{-1/z})\)).

\emph{Signs.} As \(z\to0^+\), \(F(z)\to\lambda\zeta+(1-\lambda)z_c>0\). At \(z=a\), using \(c\le1\),
\(\zeta\le\zeta_+\) and \((1-\lambda)z_c\le z_c\), \(F(a)\le\lambda\zeta_++z_c+Xe^{-1/a}-a<0\) by~\eqref{W1}. At
\(z=b\), using \(c\ge c_-\) and \(1+Ke^{-1/b}<1+K\), \(F(b)>Xc_-e^{-1/b}/(1+K)-b>0\) by~\eqref{W2}. For
\(z\ge z^*:=\lambda\zeta_++(1-\lambda)z_c+X/(1+K)\), \(F(z)<z^*-z\le0\). So \(F\) has a zero in each of \((0,a)\),
\((a,b)\) and \((b,z^*)\). By \zcref{lem:three} and \zcref{lem:zform} there are no others. This proves part~1 and the
bound on \(z_{\mathrm{out}}\).

\emph{Stability.} At \(z_L<a<1/2\), \(\psi(z_L)<\psi(a)\); at \(z_H>b>1/2\), \(\psi(z_H)<\psi(b)\). For
\(z\in\{z_L,z_H\}\), \zcref{lem:zform} and~\eqref{W3} give
\[
F'(z)+1=Xc\frac{\psi(z)}{(1+k)^2}\le X\psi(z)<\frac1{1+K}<\frac1{1+k}.
\]
Hence \(F'(z)<-k/(1+k)<\lambda-k/(1+k)\) and \(F'(z)<0\). By \zcref{lem:jac}(2) (with \(\mathcal E=k\) and
\(\varphi'=F'\)) the block is Hurwitz. At \(z_M\), \(F'(z_M)>0\) by \zcref{lem:three}, and \zcref{lem:jac}(1)
applies.

\emph{Outlet concentration.} \(c_{\mathrm{out}}=c/(1+k)\) with \(0<k<K\).
\end{proof}

\begin{theorem}[note={cascade design}]\label{thm:family}
Let \(n\ge1\) and \(\eps>0\). Suppose all \(n\) reactors have the same \(\lambda\in(0,1)\), \(K>0\), \(z_c>0\) (and
\(\Delta>0\) is common). Suppose that \eqref{W1}--\eqref{W3} hold with
\[
c_-=(1+K)^{-(n-1)},\qquad \zeta_+\ge\max\Bigl\{\eps,\ \frac{z_c+X/(1+K)}{1-\lambda}\Bigr\}.
\]
Then~\eqref{eq:model} has exactly \(3^n\) steady states, labelled by the words \(w\in\{L,M,H\}^n\), where the
letter \(w_k\) says in which window of \zcref{lem:window} the temperature of reactor \(k\) lies. All are hyperbolic. The
\(2^n\) states with \(w\in\{L,H\}^n\) are asymptotically stable. The other \(3^n-2^n\) states have a positive real
eigenvalue and are unstable. In particular this holds for the parameters of \zcref{thm:main}(2).
\end{theorem}
\begin{proof}
\emph{Induction.} We show, for \(k=0,\dots,n\), that the steady states of the subsystem formed by reactors
\(1,\dots,k\) are in bijection with \(\{L,M,H\}^k\), and that every outlet satisfies \((1+K)^{-k}\le c_k\le1\) and
\(0<z_k\le\zeta_+\). For \(k=0\) the only ``outlet'' is the feed \((1,\eps)\), and \(\eps\le\zeta_+\). Suppose the claim
holds for some \(k<n\). Every inlet \((c_k,z_k)\) of reactor \(k+1\) has \(c_k\ge(1+K)^{-k}\ge c_-\) and
\(0<z_k\le\zeta_+\). By \zcref{lem:window} reactor \(k+1\) has exactly three steady states, one in each window, so
each word of length \(k\) extends in exactly three ways. The new outlets satisfy \(c_{k+1}\ge c_k/(1+K)\ge(1+K)^{-k-1}\)
and
\[
z_{k+1}<\lambda\zeta_++(1-\lambda)z_c+\frac X{1+K}\le\lambda\zeta_++(1-\lambda)\zeta_+=\zeta_+ .
\]
Distinct words give distinct states: at the first letter where two words differ, the temperatures of that reactor lie
in disjoint windows. By \zcref{lem:reduction} every steady state arises in this way.

\emph{Stability.} By \zcref{lem:jac} the spectrum of the Jacobian is the union of the spectra of the blocks
\(J_k\). By \zcref{lem:window}(2), every block is Hurwitz when \(w_k\in\{L,H\}\), and has real eigenvalues of opposite
signs when \(w_k=M\). So all \(3^n\) states are hyperbolic. If \(w\in\{L,H\}^n\), the state is a hyperbolic sink,
hence asymptotically stable~\autocite[§2.9]{Perko2001}. Otherwise it has a positive real eigenvalue and is unstable.

\emph{The explicit family.} The parameters of \zcref{thm:main}(2) give, through \zcref{lem:zform},
\[
K=\frac1{4n},\quad \lambda=\frac1{200m},\quad z_c=\frac1{20},\quad \Delta=6400\,n\,m,\quad X=\lambda\Delta K=8 .
\]
(For instance \(z_c=\eps+\eps^2\eta_k=\eps+(1-20\eps)/20\).) Take \(a=3/20\), \(b=3\) and \(\zeta_+=\max\{\eps,9\}\).
Then \(\lambda\le1/200\), \(1+K\le5/4\), and \(\lambda\zeta_+=\max\{\eps,9\}/(200\max\{1,\eps\})\le9/200\). Also
\((1+K)^n\le e^{1/4}\), because \(n\log(1+\frac1{4n})\le\frac14\). The hypotheses then follow from five numerical
facts, independent of \(n\) and \(\eps\) (\texttt{check\_construction\_inequalities.py}, ball arithmetic):
\begin{align*}
&\zeta_+\ \ge\ 9\ \ge\ \frac{1/20+8}{1-1/200}=8.0905\ldots\ \ge\ \frac{z_c+X/(1+K)}{1-\lambda},\\
&\text{(W1)}\quad \lambda\zeta_++z_c+Xe^{-1/a}\le\tfrac9{200}+\tfrac1{20}+8e^{-20/3}=0.10518\ldots<\tfrac3{20},\\
&\text{(W2)}\quad Xc_-e^{-1/b}\ge\frac{8e^{-1/3}}{(1+K)^{n-1}},\\
&\phantom{\text{(W2)}}\quad 8e^{-1/3}=5.7322\ldots>3e^{1/4}=3.8520\ldots\ge(1+K)^n\,b,\\
&\text{(W3)}\quad (1+K)X\psi(\tfrac3{20})\le\tfrac54\cdot8\cdot\tfrac{400}9e^{-20/3}=0.5656\ldots<1,\\
&\phantom{\text{(W3)}}\quad (1+K)X\psi(3)\le\tfrac54\cdot\tfrac89e^{-1/3}=0.7961\ldots<1.
\end{align*}
For (W2), multiply the second inequality by \((1+K)^{-(n-1)}\) and combine it with the first.
\end{proof}

\begin{remark}[note={what the states look like}]
In the family of \zcref{thm:main}(2), a reactor in state \(L\) runs below \(3E/(20R)\) and converts less than
\(Ke^{-20/3}<10^{-3}/n\) of its inlet reactant. In state \(M\) it runs between \(3E/(20R)\) and \(3E/R\). In state
\(H\) it runs between \(3E/R\) and \(9E/R\). In every state it converts less than \(K/(1+K)=1/(4n+1)\) of what it
receives, so the whole cascade converts less than \(1-e^{-1/4}\approx22\%\) of the feed. For \(n=3\), \(\eps=1/20\),
the certificate gives, for example, \(\theta_1\approx6.6\times10^{-6}\) (\(T\approx T_0\)) in state \(L\) and
\(\theta_1\approx2557.5\) (\(T=T_0(1+\eps\theta_1)\approx129\,T_0\)) in state \(H\). These parameters are a mathematical construction, far
from any practical operating range. Two features carry it. First, the Arrhenius factor is bounded
(\(k(T)\le k_0\)), so a small \(k_0V/F\) caps the conversion of an ignited reactor; the S-shape of the heat release
then comes from the inflection of \(e^{-E/RT}\), which lies at \(T=E/(2R)\). Second, strong cooling (\(\lambda\) small) makes each
reactor almost insensitive to the temperature of its feed.
\end{remark}

\begin{remark}[note={the exponential limit \(\eps=0\)}]
The construction needs \(\eps>0\). With \(E(\theta)=e^\theta\), an ignited reactor leaves unconverted only the
fraction \(1/(1+D_ke^{\theta_k})\) of its feed, which is not bounded below as \(\theta_k\) grows; on the upper branch
it is even below one half (\(g_k>1/2\) where \(g_k''<0\)). So downstream reactors receive very different amounts of
reactant depending on which upstream reactors are ignited. \zcref{thm:upper} holds for \(\eps=0\). We do not know whether \(3^n\) is attained for
\(\eps=0\) and \(n\ge3\). The thesis's numerical example for \(n=2\) has \(\eps=0\), but its conjecture concerns the
model with \(\eps=RT_0/E\), which is positive.
\end{remark}

\begin{remark}[note={identical tanks}]
The family uses identical tanks. For two identical tanks with \(\eps=0\), Svoronos, Aris and
Stephanopoulos~\autocite{Svoronos1982} report at most seven steady states (for a given Damköhler number), at most three stable. For \(\eps>0\) and
\(n=2\), \zcref{thm:family} gives nine, four of them stable, already with identical tanks.
\end{remark}

\section{Computer verification}\label{sec:verif}

Environment: Fedora 44 under WSL, Python 3.12, SymPy 1.14, python-flint 0.9 (Arb ball arithmetic). From the folder of
this problem:
\begin{center}
\footnotesize
\begin{tabular}{p{4.6cm}p{8.0cm}l}
script & claims (PASS/FAIL, exit code) & time\\ \hline
\texttt{check\_model\_algebra.py} & SymPy identities: reduced variables (\zcref{lem:reduction,lem:zform}), \(g''\) and
\(G''\) factorisations (\zcref{lem:three}), \(\det\) and trace of \(J_k\) (\zcref{lem:jac}), block triangular structure
for \(n=3\), \(\psi'\) & 30 s\\
\texttt{check\_construction\_\allowbreak inequalities.py} & the numerical facts in the proof of \zcref{thm:family} (Arb), the
reductions on samples, thesis parameters for \(n=3,4\) & 3 s\\
\texttt{certify\_cascades.py} & for \(\eps=1/20\), \(n\le6\), and \(\eps\in\{1,5\}\), \(n\le4\): \(3^n\) steady states
isolated reactor by reactor, each verified by a Krawczyk test on~\eqref{eq:model} in the variables
\((y_k,\theta_k)\), boxes pairwise disjoint, block traces and determinants enclosed; \(2^n\) stable & 15 s\\
\texttt{certify\_unstable\_\allowbreak outer\_state.py} & \zcref{prop:unstable} & 2 s
\end{tabular}
\end{center}
The certificate for \(n=3\) and \(n=4\) with \(\eps=1/20\) lists every steady state (word, enclosures of
\(y_k,\theta_k\), box radius, traces and determinants of the diagonal blocks) in
\texttt{certificates/cascade\_n3\_eps1\_20.json} and \texttt{cascade\_n4\_eps1\_20.json}. With \(\eps=1/20\), the
thesis parameters are \(B=7\,680\,000\) (\(n=3\)) and \(10\,240\,000\) (\(n=4\)), \(\Da=e^{-20}/(4n)\), \(\tau=1\),
\(\alpha_k=1\), \(\beta_k=199\), \(\eta_k=0\). The certificates illustrate \zcref{thm:family}; its proof does not
depend on them. The Krawczyk test works on the equations of the thesis, not on the reduced variables, so it also
checks the reduction independently.

\paragraph{Limitations.} The upper bounds and the construction are proved for the model~\eqref{eq:model} as written in
the thesis: first-order kinetics, equal heat capacities, no wall or jacket dynamics. The attaining parameters are
extreme (ignited states above \(3E/R\), adiabatic rise \(6400\,n\,m\,E/R\)). Whether \(3^n\) states occur within
realistic ranges (bounded \(B\), small \(\eps\), temperatures below \(E/R\)) is not addressed, nor is the case
\(\eps=0\). Varma~\autocite{Varma1980}, Kubíček et al.~\autocite{Kubicek1980}, Retzloff et
al.~\autocite{Retzloff1992} and Li~\autocite{Li1994} were checked through abstracts and citing works only.

\paperreferences
\end{document}
